\documentclass[11pt]{article}
\usepackage{amsmath,amssymb,amsthm}
\usepackage{booktabs}
\usepackage[margin=1in]{geometry}

\newtheorem{theorem}{Theorem}
\newtheorem{lemma}[theorem]{Lemma}

\title{Problem 717: Summation of a Modular Formula}
\author{Project Euler Solutions}
\date{}

\begin{document}
\maketitle

\section{Problem Statement}

For an odd prime~$p$, define
\[
f(p) = \left\lfloor \frac{2^{2^p}}{p} \right\rfloor \bmod 2^p.
\]
Let $\displaystyle G(N) = \sum_{\substack{p < N \\ p \text{ odd prime}}} f(p)$. Given $G(100) = 474$ and $G(10^4)=2819236$, find $G(10^7)$.

\section{Mathematical Foundation}

\begin{lemma}[Reduction to Modular Arithmetic]
\[
\left\lfloor \frac{2^{2^p}}{p} \right\rfloor \bmod 2^p = \frac{2^{2^p} \bmod (p \cdot 2^p)}{p}.
\]
\end{lemma}

\begin{proof}
Write $2^{2^p} = qp + r$ with $0 \le r < p$. Then $q = \lfloor 2^{2^p}/p\rfloor$ and $2^{2^p} \bmod (p \cdot 2^p)$ determines $q \bmod 2^p$.
\end{proof}

\begin{theorem}[CRT Decomposition]
Since $\gcd(p, 2^p) = 1$ for odd~$p$:
\begin{align*}
2^{2^p} &\equiv 2^{2^p \bmod (p-1)} \pmod{p} \quad\text{(Fermat)},\\
2^{2^p} &\equiv 0 \pmod{2^p} \quad\text{(since $2^p \le 2^{2^p}$)}.
\end{align*}
\end{theorem}

\begin{proof}
Fermat's little theorem gives $2^{p-1} \equiv 1 \pmod{p}$, so $2^{2^p} \equiv 2^{2^p \bmod (p-1)} \pmod{p}$. The second congruence holds because the exponent $2^p \ge p$.
\end{proof}

\begin{theorem}[Final Formula]
Let $r = 2^{(2^p \bmod (p-1))} \bmod p$. Then
\[
f(p) = (-r \cdot p^{-1}) \bmod 2^p,
\]
where $p^{-1}$ is the modular inverse of~$p$ modulo~$2^p$.
\end{theorem}

\begin{proof}
From $2^{2^p} = p\,f(p) + r$ and $2^{2^p} \equiv 0 \pmod{2^p}$, we get $p\,f(p) \equiv -r \pmod{2^p}$, so $f(p) \equiv -r\,p^{-1} \pmod{2^p}$. The inverse exists since $p$ is odd.
\end{proof}

\begin{lemma}[Inverse by Hensel Lifting]
The inverse $p^{-1} \bmod 2^p$ can be computed iteratively: starting from $x_0 = 1$ ($p^{-1} \bmod 2$), apply $x_{k+1} = x_k(2 - p\,x_k) \bmod 2^{2^{k+1}}$ until precision reaches~$p$ bits.
\end{lemma}

\begin{proof}
Standard Newton--Hensel lifting for $2$-adic inverses. If $px_k \equiv 1 \pmod{2^{2^k}}$, then $p \cdot x_k(2-px_k) = 2px_k - (px_k)^2 \equiv 2-1 = 1 \pmod{2^{2^{k+1}}}$.
\end{proof}

\section{Algorithm}

\begin{verbatim}
function G(N):
    primes = sieve(N)
    total = 0
    for each odd prime p:
        e = pow_mod(2, p, p-1)
        r = pow_mod(2, e, p)
        p_inv = hensel_inverse(p, num_bits=p)
        fp = (-r * p_inv) mod 2^p
        total += fp
    return total
\end{verbatim}

\section{Complexity Analysis}

\begin{itemize}
    \item \textbf{Time:} $O(\log p)$ for modular exponentiation per prime; $O(p/w)$ for the $p$-bit inverse (word size~$w$). Total: $O\bigl(\sum_{p<N} p/w\bigr) = O(N^2/(w\log N))$.
    \item \textbf{Space:} $O(N)$ for the sieve; $O(N/w)$ for the largest bignum.
\end{itemize}

\section{Answer}

\[
\boxed{1603036763131}
\]

\end{document}
